\documentclass[11pt]{article}
\usepackage[margin=1in]{geometry}
\usepackage{amsmath,amssymb}
\usepackage[hidelinks]{hyperref}
\usepackage{microtype}

\newcommand{\core}{\operatorname{core}}

\title{Core-Dimension Sequences Need Not Be Eventually Recurrent\\
\large Literature-Already-Answered Identification}
\author{Run \texttt{fa\_banach\_001}, lane 2}
\date{9 August 2026}

\begin{document}
\maketitle

\noindent\textbf{Status.}
\texttt{literature\_already\_answered}. This is an exact later negative
answer, not a new counterexample produced by the run.

\section{Source question}

For a finite group $G$, a field $k$ of characteristic $p$, and a finite
dimensional $kG$-module $M$, let
\[
 c_n^G(M)=\dim_k\core(M^{\otimes n}),
\]
where the core is the direct sum of the nonprojective indecomposable summands.
Question 5.15.5 on source-PDF page 111 of Benson's memoir asks whether, for
every $M$, the numbers $c_n^G(M)$ satisfy a constant-coefficient linear
recurrence for all sufficiently large $n$ \cite{benson}. Equivalently, it asks
whether $\sum_{n\geq0}c_n^G(M)t^n$ is always rational.

The memoir notes that this is among the questions originating in the earlier
work of Benson and Symonds. Meng formulates their same general question as
the stronger Question 13.3 behind his Question 4.20 \cite{meng}.

\section{\hspace{0.35em}Later negative answer}

Meng's Theorem 4.21 (supporting-PDF page 25) constructs
$\Omega$-algebraic modules for which $c_n^G(M)$ is not eventually recursive.
The concrete example immediately following the theorem is
\[
 G=V_4,\qquad M=\Omega(k)\oplus\Omega^{-1}(k).
\]
Thus a finite group and finite-dimensional module violate the universal
claim in Benson's Question 5.15.5.

The supporting paper gives two equivalent checks. First, its asymptotic
theorem yields
\[
 c_n^G(M)\sim C\,2^n n^{1/2}
\]
on the relevant parity classes, with $C\neq0$. Coefficients of a rational
generating function are finite sums of exponential terms times polynomials
of \emph{integer} degree on residue classes; the factor $n^{1/2}$ is
incompatible with that form. Second, Example 5.3 displays the core series
\[
 (1+2z)\left(
   \frac{2}{1-4z^2}+\frac{4}{(1-4z^2)^{3/2}}
 \right),
\]
which is algebraic but not rational. Hence the answer to the source question
is no.

\section{Verification and scope}

Theorem 4.21 is stated for a wider family. Its proof uses nonzero variance of
the syzygy-index distribution; the displayed $V_4$ example has indices
$+1$ and $-1$ with equal weights and therefore has variance one. The concrete
counterexample is consequently within the proof's noninteger-exponent case,
independently of any zero-variance edge case in the broad theorem wording.

This identification settles only Question 5.15.5. It does not answer the
memoir's adjacent questions about
$\gamma_G(M\otimes M^*)$, symmetry of the completed representation ring,
the possible values of $\gamma_G(M)<2$, or algebraic integrality.

A bounded search used the exact memoir title, the phrases ``linear recurrence
relation with constant coefficients'' and ``eventually recursive'', and the
Benson--Symonds invariant/core-series terminology. It found Meng's 2026
paper, which explicitly labels the relevant Benson--Symonds question and
states that Theorem 4.21 answers it negatively.

\begin{thebibliography}{9}
\bibitem{benson}
D.~J. Benson,
\emph{Modular representation theory and commutative Banach algebras},
Memoirs of the American Mathematical Society 298 (2024), no.~1488;
\href{https://arxiv.org/abs/2008.13155}{arXiv:2008.13155}.

\bibitem{meng}
C.~Meng,
\emph{Asymptotic behavior of modular representations over abelian
$p$-groups},
\href{https://arxiv.org/abs/2603.11592}{arXiv:2603.11592}, 2026.
\end{thebibliography}

\end{document}
